\section{Objectives and Patient-Facing Endpoints}

\subsection{What the study is trying to find out}

\draftinstr{Roughly 1,100 characters. Adapt the objectives from
\texttt{inputs/phase-1-trial-protocol.zip},
\texttt{sections/sec-03-objectives.tex}, and restate each as a question rather
than a noun phrase. The primary objective asks whether the combination is safe
and feasible; the secondary objectives ask how the operation and the drug
performed; the exploratory objective asks what the advisory layer contributed.
Cite \texttt{onpremwhippl}.}

\subsection{Every endpoint, and the sentence it will let someone say}

\draftinstr{Roughly 900 characters introducing Figure 10 and the third rank it
adds. The parent protocol stops at the endpoint. A participant who reads only
the endpoint learns that the study measures an ISGPS grade B/C postoperative
pancreatic fistula rate and is no better informed. Rank three restates the
endpoint as the sentence it authorises. Cite \texttt{Bassi2017} for the fistula
grading and \texttt{Conroy2018} for the adjuvant survival context.}

\begin{pafig}[mermaid-type: flowchart TD, three ranks]
\node[mmgoal] (a) at (0,0) {Figure 10 slot\\objective to plain meaning};
\node[mmin] (b) at (0,-2.0) {Source:\\mermaid/fig-10-endpoint-meaning.md};
\draw[mmedge] (a) -- (b);
\end{pafig}
\vspace{-0.7cm}
\figcaption{Figure 10. Each objective carried through its endpoints to the plain sentence
the endpoint will authorise, with the three sentences that patients in the surveyed\\
literature ranked highest, clear margins, being alive, and no return, marked apart.}

\draftinstr{Replace the Figure 10 slot with the mermaid-type three-rank
flowchart from \texttt{mermaid/fig-10-endpoint-meaning.md}. Ranks at
$y = 0, -3.0, -6.4$; rank 2 text width 26 mm and rank 3 text width 30 mm; fan
the objective-to-endpoint edges with looseness 0.7 rather than drawing straight
diagonals through sibling boxes.}

\subsection{The endpoint table, with a column the protocol does not carry}

\draftinstr{Populate every row from
\texttt{inputs/phase-1-trial-protocol.zip}, \texttt{sections/sec-03-objectives.tex}.
The final column is the contribution of this paper and must be written for the
participant, not for the statistician. Verify the widths sum to the body
measure.}

\vspace{0.30em}

\noindent\begin{tabularx}{\textwidth}{L{2.5cm} L{4.5cm} C{1.5cm} Y}
\toprule
\textbf{Class} & \textbf{Endpoint} & \textbf{Timepoint} & \textbf{What it means for you} \\
\midrule
Primary & Dose-limiting toxicity rate, DL1 to DL3 & Day 30 & The dose you were given was tolerated at your level \\
Primary & Device- and procedure-related serious adverse event rate & Day 30 & Nothing the device did caused you serious harm \\
Primary & Clavien-Dindo grade III or higher rate & Day 30 & Your recovery needed no unplanned major intervention \\
Secondary & R0 resection rate & Day 90 & The margins came back clear \\
Secondary & ISGPS grade B/C fistula rate & Day 30 & Your pancreas join healed without a draining leak \\
Secondary & 90-day mortality & Day 90 & You were alive at three months \\
Secondary & Progression-free and overall survival & 24 months & The cancer had not returned by this visit \\
Exploratory & Advisory concordance and override rate & Per procedure & Your surgeon disagreed with the model here, and was free to \\
\bottomrule
\end{tabularx}

\vspace{0.30em}

\subsection{The endpoint registry as a typed record}

\draftinstr{Roughly 700 characters introducing Figure 11. An endpoint is a typed
record with a definition, a timepoint, an analysis population, a provenance
class, and, in this paper, a plain meaning. Drawing it as a schema forces every
field to be filled and exposes any endpoint where one is empty.}

\begin{pafig}[d2-type: sql\_table]
\node[d2key] (a) at (0,0) {Figure 11 slot\\endpoint registry};
\node[d2box] (b) at (0,-1.8) {Source:\\d2/fig-11-endpoint-registry.d2};
\draw[d2edge] (a) -- (b);
\end{pafig}
\vspace{-0.7cm}
\figcaption{Figure 11. The endpoint registry drawn as typed records with foreign keys to
the analysis populations and to the provenance classes, and with a plain-meaning row\\
that no clinical protocol carries but that this paper treats as a required field.}

\draftinstr{Replace the Figure 11 slot with the D2-type SQL-table figure from
\texttt{d2/fig-11-endpoint-registry.d2}. Five \texttt{d2sql} nodes at the
coordinates given in that file; each row a 4.4 mm ruled band; the plain-meaning
row filled \texttt{pablue2} to mark it as the added field; relations drawn with
crow-foot ends, the two cross-rank relations at looseness 0.85.}

\subsection{What a positive result would and would not establish}

\draftinstr{Roughly 900 characters, and write it as a proponent who intends to
be believed. A clean safety readout at day 30 and an R0 rate consistent with the
nationwide comparator would establish that the combination can proceed to a
study that could test survival. It would not establish that the combination
improves survival, and this paper does not claim otherwise. Cite
\texttt{DutchCohort2025}, \texttt{FDARobot2022}, and \texttt{NCITrial2024}.}
